\label{ch:modular}

\section{Inclusi\'on hexada--octada y realizaciones funcionales asociadas}

Fijemos un conjunto \(H\) de seis puntos y un conjunto \(O\) de ocho
puntos con una inclusi\'on distinguida \(j:H\hookrightarrow O\). En la
residencia de Witt, \(H\) es una hexada de un dodecad y \(O\) su octada
elevada; en este cap\'itulo s\'olo se usa la inclusi\'on ya fijada, no el
valor terminal de Barbero--Immirzi. Escribimos
\begin{equation}
\label{eq:modular-pairs-H}
\binom H2:=\{P\subset H:|P|=2\},
\qquad \left|\binom H2\right|=\binom62=15,
\end{equation}
y definimos dos espacios de incidencias con tipos distintos:
\begin{equation}
\label{eq:modular-F6-F8}
\mathcal F_6:=H\times\binom H2,
\qquad
\mathcal F_8:=O\times\binom H2.
\end{equation}
La inclusi\'on de incidencias es
\begin{equation}
\label{eq:modular-iota68}
\iota_{6,8}:\mathcal F_6\hookrightarrow\mathcal F_8,
\qquad
(h,P)\longmapsto(j(h),P).
\end{equation}
Es inyectiva porque \(j\) lo es y conserva el par marcado \(P\).

\begin{theorem}[Conteos de incidencia y defecto orientado]
\label{thm:modular-incidence-90-120}
Los dos espacios de \cref{eq:modular-F6-F8} satisfacen
\begin{equation}
\label{eq:modular-counts-90-120}
|\mathcal F_6|=6\binom62=90,
\qquad
|\mathcal F_8|=8\binom62=120,
\qquad
|\mathcal F_8\setminus\iota_{6,8}(\mathcal F_6)|=30.
\end{equation}
La normalizaci\'on por las veinticuatro coordenadas del par
dodecad--veinticuatro y por los quince pares de \(H\) produce
\begin{equation}
\label{eq:modular-incidence-defect}
\boxed{
\delta_{6,8}^{\rm inc}
=\frac{|\mathcal F_6|-|\mathcal F_8|}
{24\binom62}=-\frac1{12}.}
\end{equation}
La misma fracci\'on se obtiene mediante el funcional rec\'iproco asociado
a la memoria incorporada,
\begin{equation}
\label{eq:modular-reciprocal-defect}
\boxed{
\delta_{6,8}^{\rm rec}
=\frac{30}{120}-\frac{30}{90}=-\frac1{12}.}
\end{equation}
Son dos realizaciones exactas del mismo defecto orientado, pero tienen
dominios distintos y no se identifican con \(\zeta(-1)\) sin un
entrelazador adicional.
\end{theorem}

\begin{proof}
El principio multiplicativo aplicado a
\cref{eq:modular-F6-F8} da \(6\cdot15=90\) y \(8\cdot15=120\).
La imagen de \(\iota_{6,8}\) tiene noventa elementos, de modo que el
complemento tiene treinta. Sustituir en
\cref{eq:modular-incidence-defect} da
\(-30/(24\cdot15)=-1/12\), mientras
\cref{eq:modular-reciprocal-defect} es
\(1/4-1/3=-1/12\).
\end{proof}

La transici\'on hexada--octada admite dos realizaciones funcionales que no
se deducen mutuamente. El diagrama algebraico, formulado sobre los espacios
que justifican los conteos, es
\begin{equation}
\label{eq:modular-68-two-readers}
\begin{aligned}
(H\xhookrightarrow{\,j\,}O)
&\longmapsto
\left(|\mathcal F_6|,|\mathcal F_8|\right)=(90,120)
\longmapsto\Gamma_{6\to8}=\gammaH,\\
(h_6,o_8,t)=(6,8,3)
&\longmapsto-\frac{o_8-h_6}{t}=-\frac23,
\end{aligned}
\end{equation}
donde \(-2/3\) es el coeficiente de la tercera coordenada de la fila
\(13\) del constructor CKM. La primera realizaci\'on transforma la
incidencia en \`area y memoria modular; la segunda transforma la misma
transici\'on en par\'ametros de mezcla angular.
No se postula \(\boldsymbol\theta_{\rm CKM}=f(\gammaH)\): ambas salidas
conservan un antecedente com\'un.

El morfismo que determina los sectores \(90/120\) es
\(\iota_{6,8}\). La magnitud \(\gammaH\) ya ha sido determinada internamente
por la construcci\'on hexada--octada y se reutiliza aqu\'i sin reconstruirla a
partir de su representaci\'on decimal. Las
consecuencias operatorias comienzan s\'olo despu\'es de fijar la escala
elemental \(a_0\), que se define a continuaci\'on sin usar entrop\'ia ni el
producto terminal \(\gammaH a_0\).

\section{Correspondencia entre el interior non\'adico y los grados de libertad de borde}

Sea \(\mathbb T_3=\{0,1,2\}\). Todo \(x\in\mathbb Z/9\mathbb Z\) posee
una \`unica escritura de d\'igitos
\(x=x_0+3x_1\), con \(x_i\in\mathbb T_3\), y todo
\(y\in\mathbb Z/27\mathbb Z\) una \`unica escritura
\(y=y_0+3y_1+9y_2\). Por ello
\begin{equation}
\label{eq:modular-bulk-boundary-729}
\mathcal B=(\mathbb Z/9\mathbb Z)^3,
\qquad
\mathcal B_\partial=(\mathbb Z/27\mathbb Z)^2,
\qquad
|\mathcal B|=|\mathcal B_\partial|=3^6=729.
\end{equation}
La igualdad de cardinales se acompa\~na del mapa expl\'icito. Si
\(x_j=r_{2j-1}+3r_{2j}\) para \(j=1,2,3\), definimos
\begin{equation}
\label{eq:modular-beta729}
\beta_{729}(x_1,x_2,x_3)
=\bigl(r_1+3r_2+9r_3,\;r_4+3r_5+9r_6\bigr).
\end{equation}

\begin{proposition}[Bijecci\'on geneal\'ogica \(729\leftrightarrow729\)]
\label{prop:modular-beta729}
El mapa \(\beta_{729}:\mathcal B\to\mathcal B_\partial\) es biyectivo y
preserva la palabra ordenada de seis trits. No es, en general, un
isomorfismo de grupos aditivos: reagrupar los d\'igitos cambia d\'onde se
registra el acarreo.
\end{proposition}

\begin{proof}
Extraer los seis d\'igitos de tres coordenadas non\'adicas y agruparlos en
dos ternas es una permutaci\'on de coordenadas de \(\mathbb T_3^6\). La
descomposici\'on posicional es \`unica en ambos lados, por lo que la
operaci\'on tiene inversa. La advertencia aditiva se comprueba, por ejemplo,
al sumar palabras que generan acarreo entre \(r_3\) y \(r_4\): esa
frontera separa las dos coordenadas de orden veintisiete, pero atraviesa
dos coordenadas distintas en el empaquetado non\'adico.
\end{proof}

En el toro de borde \(\mathcal B_\partial\), tomamos el bloque de v\'ertices
\begin{equation}
\label{eq:modular-block-A}
A=\{0,1,\ldots,8\}\times\{0,1,\ldots,8\}.
\end{equation}
Sus enlaces orientados que cruzan el borde se separan por direcci\'on:
\begin{align}
\partial A_{90}
&=\{((-1,j),(0,j)),((8,j),(9,j)):0\leq j\leq8\},
\label{eq:modular-boundary90}\\
\partial A_{120}
&=\{((i,-1),(i,0)),((i,8),(i,9)):0\leq i\leq8\},
\label{eq:modular-boundary120}
\end{align}
donde todas las coordenadas se leen m\'odulo veintisiete. As\'i
\begin{equation}
\label{eq:modular-boundary-counts}
\partial A=\partial A_{90}\sqcup\partial A_{120},
\qquad
|\partial A_{90}|=|\partial A_{120}|=18.
\end{equation}
Los sub\'indices \(90,120\) etiquetan la procedencia de canal; no dicen que
el borde tenga noventa o ciento veinte enlaces.

\begin{theorem}[Ley areal del corte tri\'adico]
\label{thm:modular-triadic-cut}
En el grafo c\'ubico \(N\times N\times N\) con capacidad unitaria, todo
corte entre dos caras opuestas posee capacidad al menos \(N^2\), y una
capa transversal alcanza esa cota. Para \(N=9\),
\begin{equation}
\label{eq:modular-mincut-entropy}
\mathcal A_{\min}=81,
\qquad
S_{\rm tri}=81\log3.
\end{equation}
\end{theorem}

\begin{proof}
Hay \(N^2\) rectas de aristas paralelas a la direcci\'on normal, disjuntas
dos a dos. Todo corte debe interceptar cada una, de modo que su capacidad
es al menos \(N^2\). Una sola capa transversal contiene exactamente
\(N^2\) aristas y realiza el m\'inimo. Si cada enlace cortado porta un trit
m\'aximamente mezclado, su entrop\'ia es \(\log3\); la aditividad sobre los
ochenta y un enlaces da \cref{eq:modular-mincut-entropy}.
\end{proof}

La entrop\'ia \(S_{\rm tri}\) mide capacidad del corte c\'ubico. La
entrop\'ia modular de los treinta y seis enlaces de
\cref{eq:modular-boundary-counts} mide otro estado y no debe igualarse con
ella.

\section{Entrelazador colectivo de incidencia entre interior y borde}

En la suma
\(\ell^2(\mathcal F_6)\oplus\ell^2(\mathcal F_8)\), definimos
\begin{equation}
\label{eq:modular-u90-u120}
u_{90}=\frac1{\sqrt{90}}\sum_{f\in\mathcal F_6}(\delta_f,0),
\qquad
u_{120}=\frac1{\sqrt{120}}\sum_{f\in\mathcal F_8}(0,\delta_f).
\end{equation}
En
\(\ell^2(\partial A_{90})\oplus\ell^2(\partial A_{120})\), sean
\begin{equation}
\label{eq:modular-v90-v120}
v_{90}=\frac1{\sqrt{18}}\sum_{e\in\partial A_{90}}(\delta_e,0),
\qquad
v_{120}=\frac1{\sqrt{18}}\sum_{e\in\partial A_{120}}(0,\delta_e).
\end{equation}
Las parejas \((u_{90},u_{120})\) y \((v_{90},v_{120})\) son ortonormales.
Por tanto existe una \`unica isometr\'ia entre los subespacios colectivos
que satisface
\begin{equation}
\label{eq:modular-J68}
\mathcal J_{6\to8}u_{90}=v_{90},
\qquad
\mathcal J_{6\to8}u_{120}=v_{120}.
\end{equation}
Definamos
\begin{align}
D_{\rm inc}
&=90|u_{90}\rangle\langle u_{90}|
+120|u_{120}\rangle\langle u_{120}|,
\label{eq:modular-Dinc}\\
D_{\rm area}
&=90|v_{90}\rangle\langle v_{90}|
+120|v_{120}\rangle\langle v_{120}|.
\label{eq:modular-Darea}
\end{align}

\begin{theorem}[Transporte exacto de las dos etiquetas]
\label{thm:modular-J68}
La aplicaci\'on \(\mathcal J_{6\to8}\) es unitaria entre los dos planos
colectivos y entrelaza los operadores de incidencia y \`area:
\begin{equation}
\label{eq:modular-J-intertwining}
\boxed{
\mathcal J_{6\to8}D_{\rm inc}
=D_{\rm area}\mathcal J_{6\to8}.}
\end{equation}
Por tanto, toda combinaci\'on polin\'omica de las etiquetas \(90,120\), en
particular la combinaci\'on primitiva \(12:-1\), se transporta con los
mismos coeficientes.
\end{theorem}

\begin{proof}
La ortonormalidad muestra que \(\mathcal J_{6\to8}\) env\'ia una base
ortonormal en otra. En \(u_{90}\) ambos lados de
\cref{eq:modular-J-intertwining} valen \(90v_{90}\); en \(u_{120}\),
ambos valen \(120v_{120}\). La igualdad en una base prueba la identidad.
El c\'alculo funcional polin\'omico conserva el entrelazamiento.
\end{proof}

Este mapa no pretende inyectar noventa incidencias individuales en
dieciocho enlaces: transporta exactamente los dos modos uniformes y sus
etiquetas espectrales. Extenderlo a fluctuaciones no uniformes exigir\'ia
un nuevo mapa de incidencia; el operador modular definido aqu\'i s\'olo emplea
el subespacio colectivo que acaba de construirse.

\section{Determinación de la escala elemental del espectro del operador de área}

Sea \(C^*\) la coordenada angular conjugada, expresada en grados, y fijemos
\begin{equation}
\label{eq:modular-theta-clock-a0}
\theta_{\rm clk}=\frac{C^*}{6}\frac{\pi}{180},
\qquad
\boxed{
a_0=\frac{1+2\cos(10^\circ)}{3}\,\theta_{\rm clk}.}
\end{equation}
El factor que precede a \(\theta_{\rm clk}\) no es una correcci\'on
decimal. Es el car\'acter real normalizado del triplete de fases
\begin{equation}
\label{eq:modular-c10-character}
\frac13\Tr\operatorname{diag}
\left(1,\ee^{\ii\pi/18},\ee^{-\ii\pi/18}\right)
=\frac{1+2\cos(10^\circ)}3.
\end{equation}
As\'i, \(a_0\) es la media espectral de la estructura angular sobre el trit
orientado. Es positiva, se fija antes del operador de \`area y no se
selecciona a partir de la entrop\'ia. La identificaci\'on posterior
\(\gammaH=\Gamma_{6\to8}\) realiza la salida HMT en la conexi\'on de
Ashtekar--Barbero \cite{Barbero,Immirzi}; no reabre su derivaci\'on.

\begin{remark}[Alcance del s\'imbolo \(a_0\)]
En este cap\'itulo, \(a_0\equiv a_0^{\rm area}\) denota exclusivamente la
normalizaci\'on de \cref{eq:modular-theta-clock-a0}. No es el radio de Bohr
ni el factor de escala cosmol\'ogico inicial, aunque esos objetos se escriban
tradicionalmente con el mismo s\'imbolo. Ninguna identidad de este cap\'itulo
permite sustituir uno por otro.
\end{remark}
